% \noindent
% \textbf{Setup 1 (Automatic Evaluation with Simulated Users)}\quad


Automatic evaluation enables scalable testing and allows for consistent simulation of user behavior. We compare \system with the following baselines: (1) \textbf{RAG Chatbot}, a baseline that retrieves information from the search engine and interacts with the user through a one-question-one-answer paradigm. (2) \textbf{STORM + QA}, a baseline that uses the STORM framework~\citep{shao2024assisting} to generate a report for a given topic to provide general information. It then allows the user to ask follow-up questions and provides corresponding answers retrieved with the search engine.

\subsection{Evaluation Setup}
We use the \data dataset where each data point consists of an initial topic $t$ and goal $g$. 
We simulate the user 
with an LM (\texttt{gpt-4o-2024-05-13}) prompted with $t$, $g$, the discourse history $\mathcal{D}$, and the instruction for question generation. To ensure a fair comparison, we terminate the information-seeking session once it reaches 30 search queries for \system and both baselines. For all methods, the final report is generated using the two-stage approach of outline generation followed by section-by-section article generation, as proposed by STORM~\citep{shao2024assisting}, based on the interaction history. We evaluate the system quality by assessing the final report and the interaction history (\ie, discourse) with the automatic metrics defined in~\refsec{sec:eval_metrics}.
%\TODO{Ooops, w don't know what a user policy is.  Did I misread the front? If you are not going to explain the "instruction for question generation" why bother making a formal symbol policy, as if you are going to change it to different things.  } \yucheng{I make some notes in Section 2 task setup, where we used to define user policy $\pi$.}

\subsection{Automatic Metrics}
\label{sec:eval_metrics}
\noindent
\textbf{Report Quality}\quad
%We assess the quality of the complex information-seeking assistance system by evaluating the final report (see~\refsec{sec:task_setup}) 
We evaluate the final report on four aspects, \textit{Relevance}, \textit{Broad Coverage (Breadth)}, \textit{Depth}, and \textit{Novelty}, as indicators of the quality of the whole information-seeking process.\footnote{The same four aspects are used in human evaluation.} We employ Prometheus 2 \citep{kim2024prometheus}, a 7B evaluator LM, to score the report based on a 5-point rubric. To further quantify the diversity of the collected information, we also report the \textit{Information Diversity} as the average pairwise dissimilarity of $\mathcal{I}$,
\begin{equation}\nonumber
    1-\frac{\sum_{i,j\in\mathcal{I},i\neq j}\operatorname{cos}(\mathbf{i}, \mathbf{j})}{|\mathcal{I}|(|\mathcal{I}|-1)},
\end{equation}
where $\mathbf{i},\mathbf{j}$ are corresponding text embeddings obtained from OpenAI's \texttt{text-embedding-3-small}.%\footnote{\url{https://platform.openai.com/docs/guides/embeddings}}. 


%Since our task setup involves user interaction, we also report the percentage of user turns. Higher rubric scores with a lower percentage of user turns indicate better alignment with the user’s latent intent and minimal user effort. For \textit{Coverage} and \textit{Depth}, besides rubric gradings, we further employ deterministic metrics, such as \todo{information diversity and the number of unique cited URLs}.% For verifiability, we follow the methodology outlined in \citet{shao2024assisting}, calculating citation recall and citation precision as defined by \citet{gao-etal-2023-enabling}. Additionally, we utilize Mistral 7B-Instruct \citep{jiang2023mistral} to verify that cited passages support the generated sentences.

\noindent
\textbf{Discourse Quality}\quad
Since the discourse itself is valuable for human learning, we also evaluate the discourse trace using a 5-point rubric to grade each turn. This grading assesses \textit{Novelty}, \textit{Intent Alignment}, and \textit{No Repetition} for question-asking utterances (\ie, utterances with the intent \intent{Original Question} or \intent{Request Information}). For question-answering utterances that provide information%(\ie, utterances with the intent ``Potential Answer'' and ``Further Details'')
, we assess \textit{Consistency} and \textit{Engagement}. We also report the number of unique cited URLs in these utterances to indicate information diversity at the turn level. %For discourse quality, the evaluation is conducted at the utterance level and the final metric is the average result across all turns. 
Both the rubrics for report evaluation and utterance evaluation are included in Appendix~\ref{appendix:rubric}.